% ======================================================================
%  Chapter 7  --  Calibration Curves and an Analytical Development
%                 Length Equation
%
%  Split from the former combined chapter. Chapter 6 carries the
%  parametric study.
%
%  !! MISSING FROM YOUR .bib -- ADD BEFORE COMPILING !!
%     as3600ConcreteStructures2018   (AS 3600:2018)
%
%  CROSS-REFERENCES TO CHAPTER 6:
%     chap:parametric  sec:geomdb  sec:geomL  sec:matfc  sec:matstiff
%     sec:geomcover    sec:lcharsens  eq:lcharfit  eq:dbexp
%  OTHER CROSS-REFERENCES ASSUMED TO EXIST:
%     chap:tcd  eq:capacity  eq:solution
% ======================================================================

\chapter{Calibration Curves and an Analytical Development Length
         Equation}
\label{chap:devlength}

Chapter~\ref{chap:parametric} established which properties control the
predicted bond capacity and by how much. Those results were obtained one
property at a time about a single baseline, which is the right way to
isolate an influence but the wrong way to produce something a designer can
use. This chapter does the second thing.

The capacity is traced continuously against embedment length for each of
several bar diameters, producing a family of calibration curves. The family
is shown to collapse onto a single master curve whose functional form
follows from shear lag rather than from curve fitting, and that master curve
is then inverted to give a closed-form expression for development length in
which the yield stress of the bar appears exactly once. The chapter closes
by establishing the limit that the saturation of
Section~\ref{sec:geomL} imposes on that expression, which turns out to
govern its interpretation entirely.


% ======================================================================
\section{The calibration curve family}
\label{sec:curvefamily}

\subsection{Construction}

Four bar diameters between \SI{10}{} and \SI{20}{\milli\metre} were each run
at fifteen embedments between $L/d_b = 2$ and \num{30}, giving sixty cases.
All other properties are held at the baseline of
Table~\ref{tab:factorbase}: \SI{32}{\mega\pascal} concrete in a
\SI{150}{\milli\metre} cylinder, $l_{char} = \SI{20}{\milli\metre}$,
pull--push loading.

\begin{figure}[H]
\centering
% \includegraphics[width=0.85\linewidth]{figures/calibration_curves}
\caption{Calibration curves: predicted cracking load against normalised
         embedment length for four bar diameters
         ($f'_c = \SI{32}{\mega\pascal}$, $d_{cyl} = \SI{150}{\milli\metre}$,
         $l_{char} = \SI{20}{\milli\metre}$, pull--push).}
\label{fig:calcurves}
\end{figure}

Every curve has the same shape: a steep, nearly linear rise, a knee, and a
plateau flat to within a fraction of a per cent over the remaining two
thirds of the range. Mode~II governs at all sixty points; the yield load is
never reached, which is the subject of Section~\ref{sec:devlimit}.

\subsection{The effective length}

Section~\ref{sec:geomL} established the saturation for a single bar.
With four diameters available it can be quantified properly. Defining
$L_{95}$ and $L_{99}$ as the embedments at which the load first reaches
\SI{95}{} and \SI{99}{\percent} of its asymptotic value gives
Table~\ref{tab:efflength}.

\begin{table}[H]
\centering
\caption{Effective embedment length, read from the calibration curves.}
\label{tab:efflength}
\small
\begin{tabular*}{\linewidth}{@{\extracolsep{\fill}}ccccc@{}}
\hline
$d_b$ (mm) & $L_{95}$ (mm) & $L_{95}/d_b$ & $L_{99}$ (mm) & $L_{99}/d_b$ \\
\hline
10 & 73  & 7.31 & 98  & 9.77 \\
12 & 81  & 6.73 & 116 & 9.65 \\
16 & 94  & 5.85 & 158 & 9.85 \\
20 & 111 & 5.55 & 200 & 9.99 \\
\hline
\end{tabular*}
\end{table}

The two columns behave differently. At the \SI{95}{\percent} level neither
the absolute length nor the slenderness is constant: $L_{95}$ varies by a
factor of \num{1.5} across the diameter range and $L_{95}/d_b$ by
\num{1.3}. At the \SI{99}{\percent} level, however, $L_{99}/d_b$ takes the
values \num{9.77}, \num{9.65}, \num{9.85} and \num{9.99} --- a spread of
under \SI{4}{\percent}. The embedment at which the model is within one per
cent of saturation is therefore very close to ten bar diameters, whatever
the bar size.


% ======================================================================
\section{Normalisation and the master curve}
\label{sec:masterfit}

\subsection{Single-variable normalisations fail}

Plotting against $L/d_b$ does not collapse the family: at a given
slenderness the four curves differ by up to \SI{88}{\percent}. Plotting
against absolute $L$ is no better, at \SI{62}{\percent}. Dividing the
ordinate by $d_b^{0.9}$ --- the exponent established in
Section~\ref{sec:geomdb} --- reduces the spread to about \SI{27}{\percent}
but does not remove it. Optimising the exponent numerically does no better
than \SI{15}{\percent}.

\emph{No single-variable normalisation collapses the family.} This is the
first substantive result of the chapter: the slenderness ratio on which the
design codes are built is not a similarity variable for this problem.

The reason is that two independent length scales are present. The capacity
is set by the bar perimeter, which scales with $d_b$, while the distance
over which load transfers from bar to concrete is set by the shear-lag decay
of the governing equation, which has its own scaling. Any normalisation that
uses one variable must conflate them.

\subsection{The shear-lag form}

The shape in Figure~\ref{fig:calcurves} is the classical shear-lag result.
In an elastic bonded joint the interfacial shear decays from the loaded end
over a transfer length, so the load at which the peak stress reaches a
limiting value grows as the hyperbolic tangent of the bonded length divided
by that transfer length. Fitting
\begin{equation}
\label{eq:tanhfit}
F_u = F_\infty \tanh\!\left(\frac{L}{L_0}\right)
\end{equation}
to each diameter separately gives Table~\ref{tab:tanhfit}.

\begin{table}[H]
\centering
\caption{Parameters of \eqref{eq:tanhfit} fitted to each calibration curve.}
\label{tab:tanhfit}
\small
\begin{tabular*}{\linewidth}{@{\extracolsep{\fill}}cccccc@{}}
\hline
$d_b$ & $F_\infty$ & $L_0$ & $L_0/d_b$ & RMS error & max error \\
(mm)  & (kN)       & (mm)  &           & (\%)      & (\%) \\
\hline
10 & 10.825 & 44.77 & 4.48 & 3.38 & 9.63 \\
12 & 12.766 & 48.94 & 4.08 & 2.53 & 7.32 \\
16 & 16.619 & 60.96 & 3.81 & 3.10 & 10.68 \\
20 & 20.269 & 70.47 & 3.52 & 4.37 & 15.25 \\
\hline
\end{tabular*}
\end{table}

Both parameters follow clean power laws in the diameter,
\begin{equation}
\label{eq:tanhparams}
F_\infty = 1.343\, d_b^{0.907}\ \si{\kilo\newton},
\qquad
L_0 = 9.42\, d_b^{0.671}\ \si{\milli\metre},
\end{equation}
to within \SI{0.2}{} and \SI{2.1}{\percent} respectively, with $d_b$ in
millimetres. The capacity exponent \num{0.907} agrees with the \num{0.886}
and \num{0.898} obtained independently from the two diameter sweeps of
Section~\ref{sec:geomdb}, which is a useful consistency check since those
sweeps were run at single embedments and this one at fifteen.

\subsection{The master curve}

Substituting \eqref{eq:tanhparams} into \eqref{eq:tanhfit} and refitting all
sixty points simultaneously gives
\begin{equation}
\label{eq:mastercurve}
\boxed{\;
F_u = 1.349\, d_b^{0.905}
      \tanh\!\left(\frac{L}{9.32\, d_b^{0.676}}\right)\ \si{\kilo\newton}
\;}
\end{equation}
with $d_b$ and $L$ in millimetres. It reproduces every point in the family
with a root-mean-square error of \SI{3.4}{\percent}, a maximum error of
\SI{15.3}{\percent} at the shortest embedment of the largest bar, and
$R^2 = 0.993$.

\begin{figure}[H]
\centering
% \includegraphics[width=0.85\linewidth]{figures/master_curve}
\caption{The calibration curves of Figure~\ref{fig:calcurves} replotted as
         $F_u/F_\infty$ against $L/L_0$. The four families collapse onto
         $\tanh(L/L_0)$, reducing the spread from \SI{88}{} to
         \SI{12}{\percent}.}
\label{fig:mastercurve}
\end{figure}

In the normalised variables the four curves collapse to within
\SI{12}{\percent}, against \SI{88}{\percent} unnormalised.

\subsection{Why the form is mechanics and not curve fitting}

Two checks distinguish \eqref{eq:mastercurve} from an arbitrary two-parameter
fit.

The asymptotes recovered by the $\tanh$ fit agree with an independent
Richardson extrapolation of the same data to within \SI{0.1}{\percent}. The
functional form is therefore capturing the true asymptote rather than merely
passing near the last few points, which a form with the wrong tail behaviour
would not do.

The alternatives $F_\infty(1-\mathrm{e}^{-L/L_0})$ and
$F_\infty L/(L_0+L)$, fitted to the same data with the same number of free
parameters, give root-mean-square errors of \SI{4.7}{} and
\SI{8.8}{\percent} against \SI{3.3}{\percent} for $\tanh$. The form that
shear lag predicts is also the form that fits best.


% ======================================================================
\section{What sets the transfer length}
\label{sec:transferlength}

If $L_0$ were simply the elastic decay length of the governing equation, the
fit would carry no information the solution did not already contain. The
characteristic roots of \eqref{eq:solution} are complex,
$\lambda = \alpha \pm \mathrm{i}\beta$, and the envelope of the decay is set
by the real part. Table~\ref{tab:translength} compares $1/\alpha$ with the
fitted $L_0$.

\begin{table}[H]
\centering
\caption{Fitted transfer length against the elastic decay length of the
         governing equation.}
\label{tab:translength}
\small
\begin{tabular*}{\linewidth}{@{\extracolsep{\fill}}ccccc@{}}
\hline
$d_b$ & $\alpha = \operatorname{Re}\lambda$ & $1/\alpha$ & $L_0$ fitted
      & $L_0 \alpha$ \\
(mm)  & (1/mm)                              & (mm)       & (mm) & \\
\hline
10 & 0.05621 & 17.79 & 44.77 & 2.52 \\
12 & 0.05250 & 19.05 & 48.94 & 2.57 \\
16 & 0.04717 & 21.20 & 60.96 & 2.88 \\
20 & 0.04343 & 23.03 & 70.47 & 3.06 \\
\hline
\end{tabular*}
\end{table}

The two lengths are of the same order and both grow with diameter, but they
are not proportional: $1/\alpha$ follows $d_b^{0.37}$ where $L_0$ follows
$d_b^{0.67}$, and the ratio drifts by \SI{20}{\percent} across the range.

The discrepancy has the same origin as the failure of the single-variable
normalisation in Section~\ref{sec:masterfit}. The elastic solution supplies
one length scale, $1/\alpha$, which grows with the bar; the Area Method
supplies a second, the averaging radius $R_{avg} = 1.32\,l_{char}$, which is
a fixed \SI{26.4}{\milli\metre} and does not. The length over which the
\emph{criterion} sees load transferred combines the two, and $L_0$ is that
combined length rather than either separately.

This has a testable consequence. The exponent on $d_b$ in
\eqref{eq:tanhparams} is not a constant of the material: it should move
towards \num{0.37} as $l_{char}$ is reduced and the averaging radius ceases
to matter, and towards the opposite limit as $l_{char}$ grows. Repeating the
calibration family at two further values of $l_{char}$ would confirm or
refute it directly.


% ======================================================================
\section{The development length equation}
\label{sec:devequation}

\subsection{Grade independence of the calibration curve}

Nothing in \eqref{eq:mastercurve} depends on the grade of the bar. The
Mode~II capacity is a concrete cracking load determined by the geometry, the
moduli and the concrete strengths; the yield stress enters the problem only
through the Mode~I demand
\begin{equation}
\label{eq:demand}
F_I = A_s f_{sy} = \frac{\pi d_b^2 f_{sy}}{4}.
\end{equation}

The calibration curves are therefore already the grade-independent object,
and one family serves Grade~250, Grade~500 and Grade~600 alike. Raising the
grade does not move the curve; it raises the horizontal line that must
intersect it. This is what makes a single family of curves useful as a
design aid rather than a result specific to the steel it was computed with.

\subsection{Derivation}

The development length is the shortest embedment at which the concrete can
deliver the force required to yield the bar. Setting $F_u = F_I$ in
\eqref{eq:mastercurve} and inverting the hyperbolic tangent,
\begin{equation}
\label{eq:devlength}
\boxed{\;
L_d = L_0 \operatorname{artanh}\!\left(\frac{F_I}{F_\infty}\right)
    = 9.32\, d_b^{0.676}
      \operatorname{artanh}\!\left(
        \frac{\pi d_b^{2} f_{sy}}{4000 \times 1.349\, d_b^{0.905}}
      \right)\ \si{\milli\metre}
\;}
\end{equation}
with $d_b$ in millimetres and $f_{sy}$ in megapascals. The yield stress
appears exactly once, inside the inverse hyperbolic tangent; every other
quantity comes from the grade-independent calibration curve.

\subsection{Relation to the code form}

Written in terms of the mean bond stress the same statement takes the form
the codes use. Equilibrium of the embedded length gives
$\pi d_b L_d \tau_u = \pi d_b^2 f_{sy}/4$, so
\begin{equation}
\label{eq:codeform}
\frac{L_d}{d_b} = \frac{f_{sy}}{4 \tau_u},
\end{equation}
which is the basic anchorage length of fib Model Code 2010 as set out by
\textcite{cairnsBondAnchorageEmbedded2015}, with the code's design bond
stress replaced by the $\tau_u$ the present model predicts.
Equation~\eqref{eq:devlength} is the explicit solution of
\eqref{eq:codeform} once the dependence of $\tau_u$ on $L_d$ is taken from
the master curve.

\paragraph{Recovery of the code form at low utilisation.}
For $F_I \ll F_\infty$ the inverse hyperbolic tangent is linear in its
argument and \eqref{eq:devlength} reduces to
\begin{equation}
\label{eq:devlinear}
L_d \;\approx\; \frac{\pi L_0}{4 F_\infty}\, f_{sy} d_b^2
   \;=\; 5.43\times10^{-3}\, f_{sy}\, d_b^{1.77}\ \si{\milli\metre},
\end{equation}
linear in the yield stress exactly as
\textcite{as3600ConcreteStructures2018} and
\textcite{cairnsBondAnchorageEmbedded2015} are. \emph{The codes' linear
dependence on grade is the low-utilisation limit of
\eqref{eq:devlength}}, and the departure from it as $F_I$ approaches
$F_\infty$ is the model's statement that a bar worked close to the capacity
of its surrounding concrete requires disproportionately more length.

One difference survives this limit. The diameter exponent in
\eqref{eq:devlinear} is \num{1.77} where
\textcite{as3600ConcreteStructures2018} has \num{1.00}, so the model
predicts a stronger size effect than the code --- a direct consequence of
$F_\infty$ growing as $d_b^{0.91}$ while the demand grows as $d_b^2$. This
is not merely a modelling artefact. It is the same effect that led
\textcite{fosterDevelopmentLapLengths} to test bars from \SI{10}{} to
\SI{37}{\milli\metre} rather than assume a single $L_d/d_b$, and the same
effect that makes the anchorage of large high-grade bars the binding
constraint identified by \textcite{ngChangesAustralianStandards}.


% ======================================================================
\section{The limit imposed by the saturation capacity}
\label{sec:devlimit}

\subsection{The condition}

Equation~\eqref{eq:devlength} returns a real length only when
$F_I < F_\infty$. The inverse hyperbolic tangent diverges as its argument
approaches unity, which is the algebra expressing the fact that beyond the
effective length additional embedment adds no capacity. If the force
required to yield the bar exceeds the saturated capacity of the anchorage,
no embedment whatever will develop the bar.

\begin{table}[H]
\centering
\caption{Saturated capacity against the force required to yield the bar
         ($f_{sy} = \SI{550}{\mega\pascal}$,
          $f'_c = \SI{32}{\mega\pascal}$,
          $d_{cyl} = \SI{150}{\milli\metre}$,
          $l_{char} = \SI{20}{\milli\metre}$).}
\label{tab:devcheck}
\small
\begin{tabular*}{\linewidth}{@{\extracolsep{\fill}}ccccc@{}}
\hline
$d_b$ & $F_I$ & $F_\infty$ & $F_\infty/F_I$ & $f_{sy,\max}$ \\
(mm)  & (kN)  & (kN)       &                & (MPa) \\
\hline
10 & 43.2  & 10.84 & 0.250 & 138 \\
12 & 62.2  & 12.78 & 0.205 & 113 \\
16 & 110.6 & 16.58 & 0.150 & 83  \\
20 & 172.8 & 20.29 & 0.117 & 65  \\
\hline
\end{tabular*}
\end{table}

The condition fails in every case and by a wide margin. The last column
gives the yield stress that the saturated capacity could just develop: from
\SI{138}{\mega\pascal} for the \SI{10}{\milli\metre} bar down to
\SI{65}{\mega\pascal} for the \SI{20}{\milli\metre} bar, an order of
magnitude below any structural reinforcing steel. The margin worsens with
bar size, because the demand grows as $d_b^2$ while $F_\infty$ grows only as
$d_b^{0.91}$.

\subsection{The size of the discrepancy}

Two independent references fix the magnitude.

\textcite{fosterDevelopmentLapLengths} report that an anchorage length of
$16.2\,d_b$ develops the nominal yield strength of a Grade~600 bar, which
corresponds to a mean bond stress at failure of
\begin{equation}
\tau_u = \frac{f_{sy} d_b}{4 L_d}
       = \frac{600}{4 \times 16.2}
       = \SI{9.26}{\mega\pascal},
\end{equation}
against the \SI{2.24}{\mega\pascal} the model predicts for the baseline
specimen --- a factor of \num{4.1}.

The code gives the same answer. The basic development length of
Clause~13.1.2.2 of \textcite{as3600ConcreteStructures2018} for the baseline
bar, with $k_1 = k_2 = k_3 = 1$, is
\begin{equation}
\label{eq:as3600ld}
L_{sy.tb} = \frac{0.5 k_1 k_3 f_{sy} d_b}{k_2 \sqrt{f'_c}}
          = \frac{0.5 \times 550 \times 12}{\sqrt{32}}
          = \SI{583}{\milli\metre} = 48.6\,d_b ,
\end{equation}
which the model cannot reach at any embedment.

\subsection{Calibration cannot close the gap}

The obvious remedy is to calibrate $l_{char}$ upwards, since
Section~\ref{sec:lcharsens} showed it to be the only uncalibrated parameter
capable of moving the capacity substantially. The capacity must rise by a
factor of \num{4.9} for the baseline bar to be developable. Extrapolating
\eqref{eq:lcharfit}, that requires
\begin{equation}
\label{eq:lcharneeded}
l_{char} = 20 \times 4.90^{1/0.738} \approx \SI{172}{\milli\metre},
\end{equation}
which exceeds the cover of the specimen (\SI{69}{\milli\metre}) by a factor
of two and a half, and is an order of magnitude beyond any critical distance
defensible as a material length scale.

\emph{No admissible calibration of $l_{char}$ turns the model into an
ultimate-capacity predictor.} What calibration can do is fit it to measured
\emph{cracking} loads, and that is the claim Chapter~[X] should make for it.

\subsection{Why the gap exists}

Three explanations contribute, and they are not mutually exclusive.

\paragraph{Real anchorages do not saturate, because they are not elastic.}
In a physical pull-out the shear near the loaded end does not grow without
limit. Once the local bond capacity is reached there, the interface softens
and load redistributes along the bar, bringing material further from the
loaded end into play. This is precisely what the local bond--slip
formulations of \textcite{tangModelingLocalBond2020} and
\textcite{chuNewBondModel2019} represent and what the present model, being
elastic to failure, cannot. It is why codes can specify development lengths
of forty or fifty diameters at all, and it is the principal limitation of
the formulation.

\paragraph{What is predicted is first cracking, not ultimate capacity.}
The interaction index reaches unity when the averaged stresses first satisfy
the failure criterion, which corresponds to the onset of longitudinal
splitting rather than to the peak of the load--slip curve. That the
predicted mean bond stress sits just below the concrete tensile strength
--- \SI{2.24}{} against $f_t = \SI{2.83}{\mega\pascal}$ --- is consistent
with a first-cracking criterion and not with an ultimate one. The factor of
\num{4.1} is therefore largely the expected difference between cracking and
ultimate.

\paragraph{The specimen is unconfined.}
The axisymmetric idealisation has no counterpart for transverse
reinforcement, which in \textcite{408R03BondDevelopment} contributes a term
scaling with $(f'_c)^{3/4}$ and in practice is what permits the shorter
anchorages that codes allow. A model that cannot represent confinement
should not be expected to reproduce the capacity of confined anchorages.

\subsection{What the equation is therefore good for}

Equation~\eqref{eq:devlength} should be read as a \emph{splitting-initiation}
anchorage expression, not as a replacement for
\eqref{eq:as3600ld}. In that reading it has three uses. It predicts the
embedment beyond which no further capacity against splitting is gained,
which Table~\ref{tab:efflength} puts at about ten diameters. It supplies the
grade-independent shape of the capacity--embedment relationship, which the
codes assume to be linear and which it shows is not. And once $l_{char}$ is
calibrated against measured cracking loads it gives a quantitative
serviceability criterion, which is the limit state the elastic model is
actually describing.


% ======================================================================
\section{Comparison with code and research expressions}
\label{sec:codecompare}

Setting the absolute level aside, the \emph{form} of \eqref{eq:devlength}
can be compared with the expressions in use.

\begin{table}[H]
\centering
\caption{Structure of the present expression against three code and
         research formulations.}
\label{tab:codecompare}
\small
\begin{tabular*}{\linewidth}{@{\extracolsep{\fill}}llll@{}}
\hline
Source & on $f_{sy}$ & on $f'_c$ & on $c/d_b$ \\
\hline
\textcite{as3600ConcreteStructures2018}, Cl.~13.1.2.2
   & linear & $(f'_c)^{-1/2}$ & through $k_3$ \\
\textcite{cairnsBondAnchorageEmbedded2015} (MC2010)
   & linear & $(f'_c)^{-1/2}$ & through $\alpha_2$ \\
\textcite{408R03BondDevelopment}
   & linear & $(f'_c)^{-1/4}$ & explicit \\
Present model, \eqref{eq:devlength}
   & $\operatorname{artanh}$ & $(f'_c)^{-0.39}$ & $(c/d_b)^{-0.28}$ \\
\hline
\end{tabular*}
\end{table}

Three observations follow.

The model agrees with all three in the low-utilisation limit, as
\eqref{eq:devlinear} shows, and departs from them as the demand approaches
the saturated capacity, where it predicts a divergence the codes do not
represent. Whether that divergence is physical or an artefact of the elastic
formulation is exactly the question raised in Section~\ref{sec:devlimit};
the honest position is that it is the latter, and that the codes are right
to treat the relationship as linear over the range they cover.

Its exponent on concrete strength, \num{0.39}, lies between the half power
of \textcite{as3600ConcreteStructures2018} and the quarter power of
\textcite{408R03BondDevelopment}, and Section~\ref{sec:matstiff} explains
the position mechanically rather than empirically.

Its cover exponent, \num{0.28}, is of the right sign and order but weaker
than the empirical models, for the reason given in
Section~\ref{sec:geomcover}: the averaging radius makes the criterion blind
to cover beyond \SI{26.4}{\milli\metre}. This is the one structural
dependence the model gets qualitatively right and quantitatively wrong, and
it is a limitation of the criterion rather than of the stress field.


% ======================================================================
\section{Conclusions}
\label{sec:devconc}

\paragraph{The capacity--embedment relationship has a shear-lag form.}
A family of sixty calibration curves spanning a factor of two in bar
diameter is reproduced by the single expression \eqref{eq:mastercurve} with
a root-mean-square error of \SI{3.4}{\percent} and $R^2 = 0.993$. The
$\tanh$ form is not chosen for convenience: it recovers the true asymptote
to within \SI{0.1}{\percent} and fits better than the exponential and
hyperbolic alternatives tried against the same data.

\paragraph{Slenderness is not a similarity variable.}
No normalisation in $L/d_b$ alone, in $L$ alone, or in $L$ with the ordinate
scaled by any power of $d_b$ reduces the spread of the family below
\SI{15}{\percent}; the two-parameter normalisation of
\eqref{eq:mastercurve} reduces it to \SI{12}{\percent} from \SI{88}{}. Two
length scales are present --- the elastic decay length and the averaging
radius --- and $L_0$ is their combination, as
Table~\ref{tab:translength} shows.

\paragraph{A grade-independent development length expression follows by
inversion.}
Because the model's capacity contains no $f_{sy}$, the calibration curves
serve every grade and the yield stress enters \eqref{eq:devlength} exactly
once. In the low-utilisation limit the expression reduces to the linear form
used by \textcite{as3600ConcreteStructures2018} and
\textcite{cairnsBondAnchorageEmbedded2015}, so the codes' expression is
recovered as a special case, with one difference: the model carries a
diameter exponent of \num{1.77} against the code's \num{1.00}, predicting a
stronger size effect.

\paragraph{The expression is bounded by the saturation capacity, and the
bound binds.}
Development length exists only where $F_I < F_\infty$, and the saturated
capacity is between \SI{12}{} and \SI{25}{\percent} of the yield force
across the four diameters examined. The predicted mean bond stress is a
factor of \num{4.1} below that implied by
\textcite{fosterDevelopmentLapLengths}, and no admissible value of
$l_{char}$ closes the gap: \SI{172}{\milli\metre} would be required against
a specimen cover of \SI{69}{\milli\metre}.

\paragraph{The model should therefore be presented as a cracking model.}
The gap is largely the expected difference between first cracking and
ultimate capacity, compounded by the absence of post-cracking redistribution
and of confinement. Read as a splitting-initiation criterion,
\eqref{eq:devlength} is defensible and useful; read as a replacement for the
code expression it is not. Chapter~[X] should calibrate $l_{char}$ against
measured cracking loads rather than measured ultimate loads, and the
limitations stated here should frame that calibration rather than follow it.
